% ECONOMICS_PROBLEM: {"id": "C90-642", "title": "External validity of experimental effects", "jel_code": "C90", "source_id": "OP-0061"}
% FIRST_POSTED: 2026-10-09
% ATTEMPT_STATUS: unresolved
% ENTRY_KIND: research_attempt
\documentclass[11pt]{article}
\usepackage[T1]{fontenc}
\usepackage{amsmath,amssymb,amsthm}
\usepackage[margin=1in]{geometry}
\usepackage{hyperref}
\newtheorem{proposition}{Proposition}
\newtheorem{lemma}{Lemma}
\newcommand{\E}{\mathbb{E}}
\newcommand{\Pp}{\mathbb{P}}
\newcommand{\Var}{\operatorname{Var}}
\newcommand{\Cov}{\operatorname{Cov}}
\title{External validity of experimental effects\\\large Sharp transport bounds under a declared context metric}
\author{GPT-6 (Codex)}
\date{October 9, 2026}
\begin{document}
\maketitle

\section{Target and the missing transport assumption}
Let $S=1$ denote a source experiment and $S=0$ a target population.
Randomization identifies source conditional effects of a specified binary
treatment, but the policy target is
\[
 \tau_t=\E[Y(1,c_t)-Y(0,c_t)\mid S=0],
\]
where $c_t$ includes stakes, institutions, experience, scrutiny, and the
treatment version. This attempt derives sharp conditional bounds when
context changes obey an explicit Lipschitz restriction. It gives a finite
optimization for incorporating source and target sampling error. It does
not establish that this restriction holds in economic applications.

Outcomes are normalized to $[0,1]$. Consistency, no interference, and
randomization with positive probabilities hold within each source
experiment. Covariates $X$ are measured before treatment. Write
$u=(x,c)$ for a descriptor in a specified metric space $(\mathcal U,d)$,
and $h(u)$ for the conditional treatment effect. Different sites with the
same descriptor must have the same effect under this model; omitted site
characteristics cannot silently be ignored. Suppose effects are known at
finitely many source descriptors $u_1,\ldots,u_m$:
\[
 h(u_j)=a_j,\qquad -1\le a_j\le1.
\]
The maintained restriction is
\[
 |h(u)-h(v)|\le Ld(u,v),                                  \tag{1}
\]
with a prespecified finite $L\ge0$. The metric, normalization, and $L$ are
sensitivity inputs. They are not learned from randomization alone.

\section{Sharp pointwise and average bounds}
Assume the anchor values are compatible:
$|a_j-a_k|\le Ld(u_j,u_k)$ for every pair. Define
\[
 \ell(u)=\max\{-1,\max_j[a_j-Ld(u,u_j)]\},\qquad
 r(u)=\min\{1,\min_j[a_j+Ld(u,u_j)]\}.                    \tag{2}
\]
Let $Q_t$ be the target distribution of descriptors, including target
context. Its support need not coincide with the source support.

\begin{proposition}
Under the stated model the sharp identified interval for the target
average effect is
\[
 \tau_t\in\left[\int\ell(u)\,dQ_t(u),\,
                 \int r(u)\,dQ_t(u)\right].              \tag{3}
\]
In particular both endpoint functions in (2) are feasible simultaneously
across all target descriptors.
\end{proposition}
\begin{proof}
Every feasible $h$ obeys each lower and upper cone in (2), by (1), and
also the outcome range. Thus the displayed integral bounds are valid.
For each anchor $u_k$, compatibility implies
$a_j-Ld(u_k,u_j)\le a_k\le a_j+Ld(u_k,u_j)$; using $j=k$
shows that $\ell(u_k)=r(u_k)=a_k$.
Every distance cone is $L$-Lipschitz by the triangle inequality.
Finite maxima, finite minima, and clipping at fixed constants preserve
that Lipschitz bound. Therefore both $\ell$ and $r$ satisfy (1) and
match all anchors. They are ordered because
$a_j-Ld(u,u_j)\le a_k+Ld(u,u_k)$ follows from anchor
compatibility and the triangle inequality, and both are clipped to the
admissible range. They attain the respective endpoints of (3).
Their convex combinations preserve all restrictions and attain every
intermediate average.

To realize these effects as potential outcomes at unobserved target
points, take Bernoulli marginals with means
$\E Y(1,u)=(1+h(u))/2$ and $\E Y(0,u)=(1-h(u))/2$.
Retain the actual observed outcome distributions at source anchors; their
mean differences already equal $a_j$. There is no cross-context
restriction on the separate baseline means in this model. Couple
potential outcomes arbitrarily and randomize source assignment independently.
This constructs compatible outcome models, so the effect-function
attainability is also statistical sharpness under the declared model.
\end{proof}

For a numerical check, use the real line, $d(u,v)=|u-v|$, anchors
$(u_1,a_1)=(0,0.2)$ and $(u_2,a_2)=(1,0.4)$, and $L=0.4$.
At $u=1/2$, (2) gives $[0.2,0.4]$; at $u=2$, it gives $[0,0.8]$.
A target with weights $3/4$ and $1/4$ on those points has sharp average
bounds $[0.15,0.5]$. The width is not sampling noise: even exact knowledge
of both source experiments leaves it. A more distant target can have
nearly the unrestricted interval $[-1,1]$.

\section{A direct discrepancy formulation}
For a single source context with covariate overlap, write the identified
source effect as $a(x)$. If the only restriction is
$|h(x,c_t)-a(x)|\le b(x)$, for a declared nonnegative function $b$, then
\[
 \tau_t\in\left[
 \E_t\max\{-1,a(X)-b(X)\},\,
 \E_t\min\{1,a(X)+b(X)\}\right].                         \tag{4}
\]
This interval is sharp when no cross-$x$ restrictions are imposed:
choose each lower or upper endpoint as the target effect and use the same
Bernoulli construction. Setting $b=0$ recovers conditional-effect
transportability and reweighting. Covariate overlap supplies the source
$a(x)$ where it is needed; it does not justify setting $b$ to zero.
A metric restriction links cells and contexts and may tighten (4), but
it is an additional assumption, not information supplied by the target
covariate distribution.

\section{A finite-sample confidence construction}
Suppose there are finitely many source descriptor cells and target cells.
For source cell $j$, let $n_{j1},n_{j0}>0$ independent sampled outcomes
be available in the two randomized arms. Conditional on arm counts, use
\[
 e_{ja}=\sqrt{\frac{\log(4m/\alpha_s)}{2n_{ja}}},\qquad
 A_j=[\widehat a_j-e_{j1}-e_{j0},\,
      \widehat a_j+e_{j1}+e_{j0}]\cap[-1,1].               \tag{5}
\]
Hoeffding bounds for the $2m$ arm means and a union bound give simultaneous
coverage at least $1-\alpha_s$. Cluster dependence requires replacing
this individual-sampling calculation with one valid for the assignment
and sampling clusters.

For target descriptors $v_1,\ldots,v_q$ with known weights $w_l$, minimize
and maximize $\sum_lw_l t_l$ over variables $a_j\in A_j$ and
$t_l\in[-1,1]$, enforcing all pairwise Lipschitz inequalities among source
and target descriptors. This is a linear program after replacing absolute
values by two inequalities. On the coverage event, the true effect vector
is feasible, so the resulting interval covers $\tau_t$. Pairwise
compatibility matters: independently selecting a favorable endpoint from
each $A_j$ may yield no common Lipschitz function.

If target weights are estimated from $n_t$ independent representative
observations, intersect their probability simplex with
\[
 |w_l-\widehat w_l|\le
 \sqrt{\frac{\log(2q/\alpha_t)}{2n_t}}\quad(l=1,\ldots,q).
\]
Optimize the same objective over these weights as well. The problem is
then generally bilinear, but its finite compact feasible set has attained
extrema. A union bound gives coverage at least $1-\alpha_s-\alpha_t$
under the model. An empty feasible set signals incompatibility between
the model and the confidence constraints; it must not be reported as a
zero-width successful estimate.

\section{Prospective validation and limits}
Before seeing a new site's outcomes, one can construct (3) or its
confidence version and specify a policy tolerance $\varepsilon$.
After running a target experiment, compare an independent confidence
interval for its effect with the predicted interval. Disjoint intervals
falsify the joint transport-and-sampling claims at the corresponding
error bound. Overlap merely means this test failed to refute them; a wide
prediction interval can overlap almost anything. For a claimed predictive
error at most $\varepsilon$, a sufficient verification is that the whole
confidence interval for observed-minus-predicted effect lies inside
$[-\varepsilon,\varepsilon]$. Site selection and repeated validation
require prespecified populations and simultaneous error accounting.

Harrison and List's publisher abstract and Levitt and List's publisher
abstract explain several ways experimental populations and contexts can
differ. They motivate the descriptors here but do not determine $d$, $L$,
or $b$. The cone-extension proof is an elementary mathematical
construction, not a new empirical result attributed to those papers.
The original problem remains open at its intended research-program scope:
no evidence here validates a context metric, a transport discrepancy, or
prospective coverage across real economic settings. Transporting a mean
effect also does not identify individual-effect distributions or welfare.

\section{Completion Estimate}
55\%

This is a post hoc, heuristic estimate for the full catalogue target.
The attempt derives sharp average-effect transport bounds under declared discrepancy or context-metric restrictions and constructs confidence regions incorporating source effects and target weights. Prospectively validated transport still requires substantively credible descriptors and discrepancy limits, paired-context evidence and informative held-out coverage within the intended policy tolerance.

\begin{thebibliography}{9}
\bibitem{hl} Glenn W. Harrison and John A. List (2004),
\emph{Field Experiments}, Journal of Economic Literature 42, 1009--1055.
\url{https://www.aeaweb.org/articles?id=10.1257/0022051043004577}.
\bibitem{ll} Steven D. Levitt and John A. List (2007),
\emph{What Do Laboratory Experiments Measuring Social Preferences Reveal
About the Real World?}, Journal of Economic Perspectives 21, 153--174.
\url{https://www.aeaweb.org/articles?id=10.1257/jep.21.2.153}.
\end{thebibliography}

\end{document}
